% !TEX root = ../main.tex
\corrigesChapitre{11}{Statistiques et simulations}
\corrigesSouples

\begin{blocLivre}{cor-11-1}
\begin{corrige}{1}{D’un tableau aux indicateurs}
\demarche{Calculer les effectifs cumulés pour repérer les rangs de la médiane.}
\begin{questions}
\item $N=10$ ; les fréquences sont $\res{0{,}2;\ 0{,}3;\ 0{,}4;\ 0{,}1}$.
\item $\overline x=(3+8+3)/10=\res{1{,}4}$. Les valeurs aux rangs 5 et 6 sont 1 et 2 : médiane $\res{1{,}5}$. Étendue $\res{3}$.
\end{questions}
\end{corrige}
\end{blocLivre}
\begin{blocLivre}{cor-11-2}
\begin{corrige}{2}{Lire un diagramme en barres}
\demarche{Lire les effectifs sur le graphique, puis ordonner les observations.}
\begin{questions}
\item Les effectifs sont 2, 4, 3 et 1, donc $N=10$. $\overline x=(20+48+42+16)/10=\res{12{,}6}$.
\item Les valeurs centrales sont 12 et 12 : médiane 12. Les rangs des quartiles sont 3 et 8, donc $\res{Q_1=12,\ Q_3=14}$.
\end{questions}
\end{corrige}
\end{blocLivre}
\begin{blocLivre}{cor-11-3}
\begin{corrige}{3}{Tracer une boîte}
\demarche{Repérer les rangs des quartiles avant de placer les cinq repères de la boîte.}
\begin{questions}
\item La série a huit valeurs. Les deux valeurs centrales sont celles des rangs 4 et 5, soit 3 et 4 :
la médiane vaut $(3+4)/2=\res{3{,}5}$.
Les rangs des quartiles sont $\lceil8/4\rceil=2$ et $\lceil3\times8/4\rceil=6$.
On lit donc $\res{Q_1=2}$ et $\res{Q_3=4}$.
\item On place le minimum 1, les quartiles 2 et 4, la médiane 3,5 et le maximum 7.
Le trait de médiane se trouve à l’intérieur de la boîte ; les moustaches relient les bornes de la boîte aux extrêmes.
L’écart interquartile vaut $Q_3-Q_1=4-2=\res{2}$.
\end{questions}

\label{figlivre-r11-03-boite-cor}
\figureLivre[corrige]{r11-03-boite}

\end{corrige}
\end{blocLivre}
\begin{blocLivre}{cor-11-4}
\begin{corrige}{4}{Une fréquence observée}
\demarche{Une moyenne observée peut différer de l’espérance sans contredire le modèle.}
\begin{questions}
\item Quatre succès sur dix : moyenne $\res{0{,}4}$.
\item L’espérance vaut $0{,}3$, donc l’écart est $\res{0{,}1}$. Une fluctuation sur un échantillon ne suffit pas à invalider le modèle.
\end{questions}
\end{corrige}
\end{blocLivre}
\begin{blocLivre}{cor-11-5}
\begin{corrige}{5}{Comprendre un tirage simulé}
\demarche{Distinguer le résultat d’un appel et le nombre de succès sur plusieurs appels.}
\begin{questions}
\item Une loi de Bernoulli de paramètre $\res{1/4}$.
\item Le nombre de succès peut être $0,1,2,3,4$, donc la moyenne est dans $\res{\{0;1/4;1/2;3/4;1\}}$.
\item Non : $0{,}25$ est l’espérance, pas une valeur imposée à tout échantillon.
\end{questions}
\end{corrige}
\end{blocLivre}
\begin{blocLivre}{cor-11-6}
\begin{corrige}{6}{Un histogramme de largeurs inégales}
\demarche{Pour un histogramme, c’est l’aire qui porte la fréquence, pas la hauteur seule.}
\begin{questions}
\item L’aire d’un rectangle est largeur multipliée par hauteur.
Pour représenter une fréquence $f$, on choisit donc hauteur $=f/\text{largeur}$.
Les classes ont pour largeurs $10$, $20$, $30$, d’où les hauteurs
$0{,}2/10=\res{0{,}02}$, $0{,}5/20=\res{0{,}025}$ et $0{,}3/30=\res{0{,}01}$.
On vérifie les aires : $10(0{,}02)=0{,}2$, $20(0{,}025)=0{,}5$, $30(0{,}01)=0{,}3$.
\item La deuxième classe est la plus haute et représente la fréquence la plus importante.
Mais les hauteurs seules ne donnent pas les fréquences lorsque les largeurs diffèrent.
Prendre 0,2, 0,5 et 0,3 comme hauteurs donnerait des aires 2, 10 et 9, dont les proportions ne sont pas celles des données.
\end{questions}

\label{figlivre-r11-06-histogramme-cor}
\figureLivre[corrige]{r11-06-histogramme}

\end{corrige}
\end{blocLivre}
\begin{blocLivre}{cor-11-7}
\begin{corrige}{7}{Transformer une série}
\demarche{Un changement d’échelle multiplie la dispersion ; une translation conserve cette dispersion.}
\begin{questions}
\item $\overline Y=2\times12+5=\res{29}$ ; $\overline Z=\res{0}$.
\item La variance initiale vaut 9. Pour Y : variance 36 et écart type 6. Pour Z : variance 9 et écart type 3.
\item Le centrage retranche la moyenne à chaque valeur ; il déplace la série sans modifier sa dispersion.
\end{questions}
\end{corrige}
\end{blocLivre}
\begin{blocLivre}{cor-11-8}
\begin{corrige}{8}{Une aire sous une parabole}
\demarche{L’aire du carré unité vaut 1, donc la fréquence estime directement l’aire.}
\begin{questions}
\item Un point est uniforme dans un carré d’aire 1 ; la probabilité d’entrer dans la zone est donc l’aire de cette zone.
La fréquence observée fournit l’estimation $\res{137/400=0{,}3425}$.
\item L’erreur absolue par rapport à l’aire fournie est
$|137/400-1/3|=|(411-400)/1200|=\res{11/1200\approx0{,}009167}$.
La fraction reste exacte avant l’arrondi.
\item Après avoir tiré indépendamment \code{x=random()} et \code{y=random()}, le test est \code{y <= x*x}.
Les tirages appartiennent déjà au carré unité : le test supplémentaire porte uniquement sur la zone sous la parabole.
\end{questions}

\label{figlivre-r11-08-zone-cor}
\figureLivre[corrige]{r11-08-zone}

\end{corrige}
\end{blocLivre}
\begin{blocLivre}{cor-11-9}
\begin{corrige}{9}{Estimer le nombre pi}
\demarche{Le carré unité contient un quart du disque, ce qui explique le facteur 4.}
\begin{questions}
\item La fréquence dans le quart de disque est $786/1000=0{,}786$.
On estime donc $\pi$ par $4(0{,}786)=\res{3{,}144}$.
L’erreur absolue $|3{,}144-\pi|\approx0{,}002407346$ s’arrondit à $\res{0{,}00241}$.
\item Le disque de rayon 1 a pour aire $\pi$ ; le quart de disque a pour aire $\pi/4$.
Comme le carré d’échantillonnage a pour aire 1, sa fréquence estime $\pi/4$, d’où le facteur 4.
\item Un nouveau tirage peut donner un autre nombre de points dans la zone.
La taille identique ne supprime pas les fluctuations ; elle ne force pas une répétition exacte du résultat.
\end{questions}

\label{figlivre-r11-09-zone-cor}
\figureLivre[corrige]{r11-09-zone}

\end{corrige}
\end{blocLivre}
\begin{blocLivre}{cor-11-10}
\begin{corrige}{10}{Comparer des moyennes d'échantillons}
\demarche{Comparer chaque écart à la même tolérance avant de compter les moyennes retenues.}
\begin{questions}
\item $E(X)=0(1/2)+2(1/2)=1$ et $E(X^2)=0+4(1/2)=2$.
Ainsi $V(X)=2-1^2=1$, donc $\res{\mu=1,\ \sigma=1}$.
\item Les sommes des quatre échantillons sont 4, 8, 0 et 6.
Après division par 5, les moyennes valent $\res{0{,}8;\ 1{,}6;\ 0;\ 1{,}2}$.
\item Le seuil commun est $2/\sqrt5\approx0{,}8944$.
Les écarts à 1 sont $0{,}2$, $0{,}6$, 1 et $0{,}2$ : le troisième seul dépasse le seuil.
La proportion retenue est $3/4=\res{75\,\%}$.
Ce résultat décrit ces quatre échantillons, sans imposer ce pourcentage à d’autres tirages.
\end{questions}

\label{figlivre-r11-10-critere-cor}
\figureLivre[corrige]{r11-10-critere}

\end{corrige}
\end{blocLivre}
\begin{blocLivre}{cor-11-11}
\begin{corrige}{11}{Effet de la taille d'échantillon}
\demarche{Le seuil diminue comme l’inverse de la racine de la taille.}
\begin{questions}
\item Pour $n=16$, le seuil vaut $4/4=\res{1}$ ; pour $n=64$, $4/8=\res{1/2}$.
\item $|2{,}2-1|=1{,}2>1$ : non. $|1{,}4-1|=0{,}4\leq0{,}5$ : oui.
\item Le seuil varie comme $1/\sqrt n$. Pour le diviser par 3, il faut multiplier $n$ par $\res{9}$.
\end{questions}
\end{corrige}
\end{blocLivre}
\begin{blocLivre}{cor-11-12}
\begin{corrige}{12}{Simuler une loi à trois valeurs}
\demarche{Partager un seul tirage uniforme entre des intervalles de longueurs prescrites.}
\begin{questions}
\item $\mu=-0{,}5+0{,}6+1=\res{1{,}1}$. $E(X^2)=0{,}5+1{,}2+5=6{,}7$, donc $V(X)=\res{5{,}49}$ et $\sigma=\res{\sqrt{549}/10\approx2{,}343}$.
\item On partage $[0;1[$ aux seuils 0,5 et 0,8. Le code ci-dessous utilise le même tirage pour les deux tests.
\item Pour $n\geq1$, on moyenne une liste de $n$ simulations.
\item Répéter les expériences à chaque taille, puis étudier les écarts absolus à 1,1. Une moyenne plus grande ou plus petite que 1,1 reste possible.
\end{questions}
\begin{python}
from random import random
def tirage():
    u = random()
    if u < 0.5:
        return -1
    if u < 0.8:
        return 2
    return 5
def moyenne(n):
    return sum(tirage() for k in range(n))/n
\end{python}

Le deuxième test n’est atteint que si le premier a échoué : la valeur 2 correspond donc à $0{,}5\leq u<0{,}8$, un intervalle de longueur $0{,}3$. Un deuxième tirage aléatoire modifierait le modèle.

\label{figlivre-r11-12-partition-cor}
\figureLivre[corrige]{r11-12-partition}

\end{corrige}
\end{blocLivre}
\begin{blocLivre}{cor-11-13}
\begin{corrige}{13}{Répéter cinq cents échantillons}
\demarche{Simuler de nouveaux échantillons et distinguer le compteur de la proportion finale.}
\begin{questions}
\item $\mu=0{,}4$, $\sigma=\sqrt{0{,}24}$ et le seuil vaut $\res{2\sqrt{0{,}24}/5\approx0{,}19596}$.
\item Chaque moyenne est obtenue à partir de 25 nouveaux tirages. Le compteur est ensuite divisé par 500.
\item Non : il s’agit d’une proportion expérimentale dépendant du modèle et des tirages, pas d’une valeur universelle imposée.
\end{questions}
\begin{python}
from random import Random
from math import sqrt
rng = Random(2026)
compteur = 0
for j in range(500):
    m = sum(rng.random() < 0.4 for k in range(25))/25
    if abs(m - 0.4) <= 2*sqrt(0.24)/5:
        compteur = compteur + 1
proportion = compteur/500
\end{python}

Avec la graine 2026 et ce programme, on obtient 470 moyennes retenues sur 500, soit une proportion $0{,}94$. La graine rend cet exemple reproductible ; elle ne transforme pas 94\,\% en règle générale.

\end{corrige}
\end{blocLivre}
\begin{blocLivre}{cor-11-14}
\begin{corrige}{14}{Comparer deux groupes et leur réunion}
\demarche{Comparer les écarts à un centre commun pour mesurer ce que la moyenne masque.}
\begin{questions}
\item Les deux moyennes valent $\res{5}$. Pour A, $v_A=(1+0+0+1)/4=\res{1/2}$ ; pour B, $v_B=(16+0+0+16)/4=\res{8}$.
\item La moyenne reste 5. La somme des carrés des écarts vaut $2+32=34$, donc la variance est $\res{34/8=17/4}$.
\item La série ordonnée est $1,4,5,5,5,5,6,9$. Médiane 5 ; $Q_1$ au rang 2 vaut 4 ; $Q_3$ au rang 6 vaut 5.
\item Le groupe B est beaucoup plus dispersé, bien que sa moyenne soit identique. Les valeurs extrêmes sont invisibles si on ne rapporte que la moyenne.
\end{questions}
\label{figlivre-r11-14-groupes-cor}
\figureLivre[corrige]{r11-14-groupes}

\end{corrige}
\end{blocLivre}
\begin{blocLivre}{cor-11-15}
\begin{corrige}{15}{Fréquences des lettres d'un texte}
\demarche{Compter chaque lettre puis diviser par la longueur totale du texte.}
\begin{questions}
\item Dans BANANE : B et E ont fréquence $1/6$, A et N fréquence $1/3$. Dans ANANAS : A $1/2$, N $1/3$, S $1/6$. Les autres fréquences sont nulles.
\item Pour un texte non vide, on peut écrire \code{[texte.count(c)/len(texte) for c in alphabet]}, où \code{alphabet} est la chaîne des 26 lettres majuscules. La somme vaut 1 si tous les caractères sont des lettres de cet alphabet.
\item Il y a deux A et un E sur six positions : $P(V=1)=P(V=0)=1/2$, donc $\res{E(V)=1/2}$ et $\res{V(V)=1/4}$.
\end{questions}
\end{corrige}
\end{blocLivre}
